\subsection{Implementació del Snap}

El Snap constitueix el punt d’entrada del sistema dins de MetaMask i actua com a capa d’intercepció i presentació dels resultats. La seva funció principal no és realitzar tota l’anàlisi de manera interna, sinó capturar la transacció abans de la confirmació, preparar-ne la informació rellevant, consultar el \textit{backend} local i presentar el resultat a l’usuari dins la pròpia interfície de la cartera.

\subsubsection{Estructura general del Snap}

La implementació del Snap s’ha centralitzat principalment en el fitxer \texttt{packages/snap/src/index.tsx}, que conté tant la definició dels punts d’entrada del Snap com la lògica necessària per comunicar-se amb el \textit{backend}. Aquesta organització simplifica el prototip i permet concentrar en un únic punt la intercepció de transaccions, la preparació de les dades i la construcció de la interfície mostrada a MetaMask.

En una primera fase, el Snap només disposava d’un mètode bàsic de prova, accessible des de la DApp local mitjançant \texttt{wallet\_invokeSnap}. Aquest comportament inicial servia per validar la instal·lació del Snap, la comunicació amb la cartera i la generació de diàlegs simples dins MetaMask. Posteriorment, aquesta base es va ampliar per incorporar un flux real d’anàlisi de transaccions, primer mitjançant un mètode invocat manualment des de la DApp i després mitjançant la intercepció automàtica amb \texttt{onTransaction}. 

\subsubsection{Permisos i configuració}

Perquè el Snap pogués operar dins del flux de transaccions, va ser necessari definir al \textit{manifest} un conjunt de permisos específics. Entre els més importants hi ha \texttt{snap\_dialog}, que permet mostrar diàlegs dins MetaMask; \texttt{endowment:network-access}, que habilita l’ús de \texttt{fetch()} per comunicar-se amb el \textit{backend} local; i \texttt{endowment:transaction-insight}, que permet generar \textit{transaction insights} abans de la confirmació de la transacció. En aquesta darrera capacitat també es fa servir l’opció \texttt{allowTransactionOrigin} per poder disposar d’informació sobre l’origen de la petició. 

La configuració correcta d’aquests permisos va ser un pas crític en la implementació, ja que sense ells el Snap podia continuar mostrant diàlegs simples, però no podia ni consultar el \textit{backend} ni interceptar automàticament transaccions reals dins MetaMask.

\subsubsection{Interacció manual mitjançant \texttt{onRpcRequest}}

Abans d’activar la intercepció automàtica, es va implementar una primera via de prova a través de \texttt{onRpcRequest}. Aquest mecanisme permetia que la DApp local invoqués explícitament el Snap amb un mètode personalitzat, anomenat \texttt{analyzeTransaction}, enviant una estructura simplificada de transacció. El Snap rebia aquestes dades, les reenviava al \textit{backend} i mostrava el resultat amb un diàleg de MetaMask.

Aquesta fase va ser especialment útil per validar de manera incremental la integració del sistema. Gràcies a ella es va poder comprovar, abans d’entrar en la intercepció real, que la comunicació Snap $\rightarrow$ \textit{backend} $\rightarrow$ Ollama funcionava correctament i que la resposta podia representar-se dins de la UI limitada de MetaMask. Aquest enfocament també va facilitar el procés de depuració inicial, ja que permetia provar el \textit{backend} sense dependre encara del flux complet d’una transacció real. 

\subsubsection{Intercepció automàtica de transaccions amb \texttt{onTransaction}}

Un cop validada la connexió amb el \textit{backend}, la implementació es va ampliar amb el punt d’entrada \texttt{onTransaction}, que és el mecanisme proporcionat per MetaMask Snaps per interceptar transaccions abans que l’usuari les confirmi. Aquest \textit{handler} s’executa automàticament quan MetaMask detecta una transacció compatible i hi ha un Snap instal·lat amb els permisos adequats. 

Dins d’aquest \textit{handler}, el Snap extreu els camps principals de la transacció, entre els quals destaquen l’adreça d’origen, l’adreça de destinació, el valor enviat, la calldata, la xarxa utilitzada i l’origen de la transacció. Aquesta informació es transforma a una estructura JSON simplificada que després s’envia al \textit{backend} local per ser analitzada. El fet de treballar amb una estructura pròpia i normalitzada facilita tant el tractament posterior com la possible extensió del sistema amb nous tipus d’anàlisi.

\subsubsection{Comunicació amb el backend}

Per comunicar-se amb el \textit{backend}, el Snap utilitza una petició HTTP local mitjançant \texttt{fetch()} contra l’endpoint \texttt{http://localhost:3000/analyze}. Aquesta petició envia les dades de la transacció en format JSON i rep com a resposta una estructura amb el veredicte final, l’acció recomanada, els indicis del camp \texttt{findings}, el context disponible, la revisió complementària de la IA i l’explicació final. El Snap no implementa directament la lògica heurística ni la generació de text, sinó que delega aquestes responsabilitats al \textit{backend}, mantenint-se així com un component lleuger i centrat en la integració amb MetaMask.

Aquesta decisió respon tant a motius arquitectònics com a limitacions de l’entorn del Snap. D’una banda, l’aïllament i les restriccions de MetaMask Snaps fan poc convenient carregar-lo amb processament més pesat. De l’altra, separar la lògica d’anàlisi en un servei local independent facilita l’evolució del sistema, la reutilització de codi i la incorporació de millores futures sense haver de modificar constantment el Snap.

\subsubsection{Presentació dels \textit{insights} dins MetaMask}

Un cop rebuda la resposta del \textit{backend}, el Snap construeix la interfície que es mostra dins de MetaMask. Aquesta interfície es genera amb els components oficials del sistema de Snaps, com ara \texttt{Box}, \texttt{Text}, \texttt{Bold} i \texttt{Heading}, ja que l’entorn de MetaMask no permet utilitzar una interfície web arbitrària.

En les primeres versions del prototip, la informació presentada a l’usuari es limitava principalment al nivell de risc i a una explicació en llenguatge natural. Amb l’evolució del sistema, la resposta del \textit{backend} s’ha ampliat i el Snap pot mostrar un conjunt més ric d’elements, com ara el veredicte final, l’acció recomanada, els indicis principals, el context local detectat, la revisió complementària de la IA i la font que ha motivat el resultat.

La versió final unifica en català totes les etiquetes visibles i les explicacions noves. El Snap transforma els codis interns de risc \texttt{BAJO}, \texttt{MEDIO}, \texttt{ALTO} i \texttt{DESCONOCIDO} en \textit{Baix}, \textit{Mitjà}, \textit{Alt} i \textit{Desconegut}, i presenta \texttt{ALLOW}, \texttt{REVIEW} i \texttt{BLOCK} com \textit{Permetre}, \textit{Revisar} i \textit{Bloquejar}. Els codis originals no es modifiquen dins del contracte JSON ni a SQLite, per mantenir la compatibilitat amb l’historial, les proves automatitzades i els consumidors de l’API.

Aquesta ampliació permet que l’usuari no rebi només una alerta genèrica, sinó una visió més estructurada del motiu pel qual una transacció pot ser considerada segura, sospitosa o arriscada. Tot i això, el Snap continua mantenint una responsabilitat limitada: no reconstrueix l’anàlisi intern ni aplica les regles de risc, sinó que adapta la resposta rebuda al format visual permès per MetaMask.

La limitació dels components disponibles condiciona el disseny de la interfície, però també ajuda a centrar el missatge en la informació essencial: el nivell de risc, els motius principals i l’explicació final que ha de servir de suport a la decisió de l’usuari.
